\clearpage
\section*{September 24, 2026: Likelihood estimation and bootstrap intervals}
\textit{Agent-drafted research entry.}

Jacob asked for a version of the main estimation that a structural referee
would expect: a likelihood rather than equally weighted squared share gaps,
and standard errors. It is \path{bootstrap_notch_model.R} in
\path{tasks/estimate_notch_model}; the least-squares main specification is
unchanged and remains the reported fit.

\paragraph{The likelihood needs an unexplained share.} Treating each of the
278 recent parents of at most 300 units as a draw from the predicted cell
probabilities fails at every point of the grid: some observed parent always
has probability zero. At the least-squares estimate these are three single
buildings of 105--119 units, which the jump makes strictly worse than 99, and a
three-building parent of 180--197 units, which two buildings could hold.
Adding a share of parents who keep their 2019--2022 outcome does not help,
since the historical sample has no parents in four observed cells, $99+99$
among them. The likelihood therefore lets a share $\varepsilon$ of recorded
recent outcomes be unexplained by the model, as linkage or unit-count errors
would be, spread uniformly over sizes 50--300 and one, two or three-plus
buildings. This describes the records, not developer behavior, so units lost
still use the model's choices only. It is not the non-optimizer mixture, which
Jacob wants to compare with heterogeneity in the jump later.

\paragraph{The bootstrap.} Each of 500 draws resamples parents with
replacement within period and borough, recalibrates the weights to the
resampled recent parents, and re-estimates both estimators on the full grid.
Stratifying by borough keeps Staten Island (five historical and two recent
parents) in every calibration. The model's choices do not depend on the
weights, so one pass over the grid scores every draw. Intervals are the 2.5th
and 97.5th percentiles of the draws; likelihood-ratio intervals are the grid
values within 1.92 log points of the maximum profile likelihood.

\begin{center}
\small
\begin{tabular}{lrrrrr}
\hline
 & \multicolumn{2}{c}{Least squares} & \multicolumn{3}{c}{Likelihood} \\
 & Estimate & Bootstrap & Estimate & Bootstrap & Likelihood ratio \\
\hline
Jump $\kappa$ & 0.10 & 0.04--0.24 & 0.12 & 0.03--0.22 & 0.04--0.18 \\
Kink $\tau$ & 0.05 & 0--0.15 & 0.05 & 0--0.30 & 0.02--0.20 \\
Cost growth $\gamma$ & 1.75 & 1--3 & 1.5 & 1--3 & 1--3 \\
Mean splitting cost $\sigma$ & 0.25 & 0.10--1.0 & 0.32 & 0.08--1.0 & 0.16--0.63 \\
Unexplained share $\varepsilon$ & -- & -- & 0.03 & 0.001--0.075 & -- \\
Units lost & 1,018 & 396--2,197 & 1,130 & 337--2,981 & -- \\
\hline
\end{tabular}
\end{center}

\paragraph{What it shows.} The likelihood and least squares agree closely at
the point estimates, so the equal weighting of cells was not driving the fit.
Uncertainty is large. Units lost range over a factor of five to nine: the
estimate of about 1,000 is best read as somewhere between 400 and 2,200 under
least squares. The growth of the splitting cost is not identified: its
bootstrap draws and its likelihood-ratio interval cover the whole grid from 1
to 3. The jump is bimodal in the likelihood draws, with about 40 percent near
0.03--0.04 and most of the rest at 0.10--0.14; the least-squares draws center
on 0.10--0.12. The kink's likelihood-ratio interval excludes zero, but 9 to 12
percent of bootstrap draws choose no kink. The unexplained share is 3 percent
of recent parents, about eight; in 22 percent of draws it falls to the grid
minimum, presumably draws that happen to omit the parents in impossible cells.

\paragraph{Limits.} Percentile intervals on a grid inherit the grid's bounds,
most visibly for $\gamma$ and the upper end of $\sigma$. Stratifying by
borough holds borough counts fixed, which slightly understates uncertainty.
The uniform spread of unexplained outcomes is a convenience, not an estimate of
how records go wrong.

\paragraph{Next.} Compare the non-optimizer share with heterogeneity in the
jump. Jacob reserves the test that parents in areas with higher pre-policy
construction wages react less, because 485-x raises their costs less, for the
final version of the model.
